\lecture{2.2}{The Roadmap for Causal Analysis}
\label{lec:2-2}

\subsection{Overview}

Causal analysis is easiest to learn as a sequence of distinct tasks. Mixing these tasks is a common source of confusion. For example, fitting a regression model before defining the causal question often leads to interpreting a convenient coefficient rather than estimating a meaningful causal parameter.

\begin{keyidea}
The causal roadmap separates scientific definition, causal identification, statistical estimation, and interpretation. Each step answers a different question.
\end{keyidea}

\begin{figure}[H]
\centering
\begin{tikzpicture}[node distance=.75cm,>=Latex]
  \tikzset{box/.style={draw, rounded corners, fill=blue!4, text width=10cm, minimum height=.75cm}}
  \node[box] (s1) {1. Specify a causal model};
  \node[box, below=of s1] (s2) {2. Specify a causal question and causal parameter};
  \node[box, below=of s2] (s3) {3. Link observed data to the causal model};
  \node[box, below=of s3] (s4) {4. Assess identifiability};
  \node[box, below=of s4] (s5) {5. Specify a statistical parameter and statistical model};
  \node[box, below=of s5] (s6) {6. Estimate};
  \node[box, below=of s6] (s7) {7. Interpret};
  \foreach \a/\b in {s1/s2,s2/s3,s3/s4,s4/s5,s5/s6,s6/s7}
    \draw[->, thick] (\a) -- (\b);
\end{tikzpicture}
\caption{A practical roadmap for causal analysis.}
\end{figure}

\subsection{Specify a Causal Model}

A causal model represents scientific knowledge about how variables in the system relate to one another. At this stage we ask questions such as:
\begin{itemize}
  \item What factors affect whether a patient receives early imaging?
  \item What factors affect later disability?
  \item Which variables occur before treatment, during follow-up, or after the outcome?
  \item Which important causes were not measured?
\end{itemize}

In this course, causal models will often be represented by directed acyclic graphs (DAGs) and single world intervention graphs (SWIGs). The graph is not just a picture. It encodes assumptions about causes, noncauses, and conditional independencies.

\subsection{Specify a Causal Question and Parameter}

The causal question must state the intervention and the outcome. A vague question such as "Does imaging help?" is not yet a causal parameter. A more precise version is:
\begin{quote}
Among older adults with a new primary care visit for back pain, what is the difference in mean 12-month disability if everyone received lumbar imaging within six weeks versus if no one received lumbar imaging within six weeks?
\end{quote}

Using potential outcome notation, this question could be represented by
\[
\psi=\E\{\Yone\}-\E\{\Yzero\},
\]
where $\Yone$ is the outcome that would be observed under early imaging and $\Yzero$ is the outcome that would be observed under no early imaging.

\begin{definition}[Causal parameter]
A \emph{causal parameter} is a summary of the distribution of counterfactual variables in a target population. It is the causal estimand.
\end{definition}

\subsection{Link Observed Data to the Causal Model}

The causal model describes variables that may include both observed and unobserved quantities. The observed data describe what was actually measured. A common observed data structure for a point-treatment study is
\[
O=(W,A,Y),
\]
where $W$ denotes baseline covariates, $A$ observed treatment, and $Y$ observed outcome.

The analyst must state how these observed variables correspond to nodes in the causal model. If the graph contains an unmeasured severity variable $U$, the observed data do not include $U$. If the graph contains a broad baseline node $W$, the analyst must clarify which measured variables are included in $W$.

This step also includes assumptions that connect observed outcomes to counterfactual outcomes, especially consistency, introduced later in this chapter.

\subsection{Assess Identifiability}

Identification asks whether the causal parameter can be expressed as a function of the observed data distribution under the stated causal model.

\begin{definition}[Identifiability]
A causal parameter is \emph{identifiable} from the observed data distribution under a set of assumptions if all causal models satisfying those assumptions and producing the same observed data distribution imply the same value of the causal parameter.
\end{definition}

Plainly, identifiability asks: if we had infinite observed data from this study design, would the causal question be answerable? This question comes before estimation. No estimator, no matter how flexible, can recover a causal parameter that is not identified without additional assumptions.

\begin{identificationbox}
Identification is the bridge from counterfactual quantities to observed-data quantities. A typical goal is to replace a causal expression such as $\E\{\Ya{a}\}$ with an observed-data expression such as $\E\{\E(Y\given A=a,W)\}$ under appropriate assumptions.
\end{identificationbox}

\subsection{Specify Statistical Parameter and Model}

Once a causal parameter has been identified, the target becomes a statistical parameter of the observed data distribution. For example, an identification result might show that
\[
\E\{\Ya{a}\}=\E\{\E(Y\given A=a,W)\}.
\]
The right-hand side is an observed-data parameter. It depends on the distribution of $(W,A,Y)$, not directly on unobserved counterfactuals.

The statistical model describes what is assumed about the observed data distribution. A nonparametric model makes few restrictions. A parametric model might assume a specific logistic regression or linear regression form. The choice of statistical model affects estimation and inference, but it is conceptually separate from the causal identification assumptions.

\subsection{Estimate}

Estimation uses the finite sample to approximate the statistical parameter. If the identified parameter involves the conditional mean $\E(Y\given A=a,W)$, then an estimator may require fitting an outcome regression. If it involves $\Prob(A=a\given W)$, then an estimator may require fitting a propensity score.

\begin{estimationbox}
The estimation step answers a statistical question: given the identified observed-data parameter and a sample $O_1,\ldots,O_n$, how should we compute a numerical estimate and quantify uncertainty?
\end{estimationbox}

Estimation should not change the estimand. A regression coefficient may be easy to estimate, but it is not automatically the causal parameter. The target parameter should be defined before the estimator is chosen.

\subsection{Interpret}

The final step translates the estimate back into the scientific question while respecting the assumptions and limitations. Interpretation should state:
\begin{itemize}
  \item the target population;
  \item the intervention being compared;
  \item the outcome and time scale;
  \item the causal assumptions required;
  \item the statistical uncertainty;
  \item limitations from measurement, missing data, adherence, and generalizability.
\end{itemize}

A careful interpretation might say that the result estimates the mean difference in 12-month disability under two specified imaging policies, among patients represented by the study population, assuming the measured covariates were sufficient to control confounding, treatment was well-defined, and follow-up bias was adequately handled.
